% subsection: 階差型 \, $a_{n+1}=a_n+f(n)$

  \subsection{\texorpdfstring{階差型 \, $a_{n+1}=a_n+f(n)$}{階差型}}\label{節：階差型}
  \begin{bookpoint}
    \textbf{変化を足す}\par
    差が添字ごとに変わっても,初項からの変化を足せば元の項に戻る.和の範囲は実際の更新の回数に合わせる.
  \end{bookpoint}
    足す量が一定でなくても、初項からの変化を足し戻せる。
    一歩目では$f(1)$、二歩目では$f(2)$を足す。
    直前の項を代入すると、その変化が順に並ぶ。
    \begin{align*}
      &a_{n+1}=a_n+f(n)\\
      &a_{n+1}=\{a_{n-1}+f(n-1)\}+f(n)\\
      &a_{n+1}=\left[\{a_{n-2}+f(n-2)\}+f(n-1)\right]+f(n)\\
      &\qquad \vdots\\
      &a_{n+1}=a_1+f(1)+f(2)+\cdots+f(n)\\
      &a_{n+1}=a_1+\sum_{k=1}^{n} f(k)\\
      &\therefore \quad a_n=a_1+\sum_{k=1}^{n-1} f(k)
    \end{align*}
    変化を足すところまでで、元の\bookterm{sequence}{数列}は和の形に表せた。
    残るのは、その和の計算である。
    $f(n)$の形に合わせて、\bookterm{arithmetic}{等差数列}の和、\bookterm{geometric}{等比数列}の和、\bookterm{partial}{部分分数分解}を使う。

    \noindent \bookblocklink{example-difference-1}{problem}{answer}{【例題】}\par\nobreak
    \begin{enumerate}[label=(\arabic*),parsep=3pt,format=\bookitemlink{example-difference-1}{problem}{answer}]
      \item $\displaystyle a_1=1,\, a_{n+1}=a_n+n$
      \item $\displaystyle a_1=1,\, a_{n+1}=a_n+2^n$
      \item $\displaystyle a_1=0,\, a_{n+1}=a_n+\frac{1}{n(n+1)}$
      \item $\displaystyle a_1=1,\, a_{n+1}=a_n+(n+1)^2$
    \end{enumerate}

    \noindent\bookblocklink{example-difference-1}{hint}{problem}{【方針】}\par\nobreak
    \begin{enumerate}[label=(\arabic*),parsep=3pt,format=\bookitemlink{example-difference-1}{hint}{problem}]
      \item 足す量は$n$です。初項からの変化を、一次式の和にまとめます。
      \item もちろん,\,$f(n)$となる関数には指数関数も含まれます.何にせよやることは\,(1)と同じです.
      \item \bookterm{sigma}{シグマ}は\bookterm{partial}{部分分数分解}で処理しましょう.
      \item 足す量は$(n+1)^2$です。初項を含めて、平方の和としてまとめます。
    \end{enumerate}

    \noindent\bookblocklink{example-difference-1}{answer}{problem}{【解答】}\par\nobreak
    \begin{enumerate}[label=(\arabic*),parsep=3pt,format=\bookitemlink{example-difference-1}{answer}{problem}]
      \item $\displaystyle a_n=1+\frac{n(n-1)}{2}$
      \item $\displaystyle a_n=2^n-1$
      \item $\displaystyle a_n=1-\frac{1}{n}$
      \item $\displaystyle a_n=\frac{1}{6}n(n+1)(2n+1)$
    \end{enumerate}

    \noindent\bookblocklink{example-difference-1}{solution}{problem}{【解法】}\par\nobreak
    \begin{enumerate}[label=(\arabic*),parsep=3pt,format=\bookitemlink{example-difference-1}{solution}{problem}]
      \item $\displaystyle a_1=1,\,a_{n+1}=a_n+n$\\
            \bookterm{difference}{階差}型の\bookterm{general}{一般項}
            \[
              a_n=a_1+\sum_{k=1}^{n-1}f(k)
            \]
            を用いると,
            \begin{align*}
              a_n
                &=1+\sum_{k=1}^{n-1}k\\
                &=1+\frac{n(n-1)}{2}
            \end{align*}

      \item $\displaystyle a_1=1,\,a_{n+1}=a_n+2^n$\\
            同様に,
            \begin{align*}
              a_n
                &=1+\sum_{k=1}^{n-1}2^k\\
                &=1+(2+2^2+\cdots+2^{n-1})\\
                &=1+(2^n-2)\\
                &=2^n-1
            \end{align*}

      \item $\displaystyle a_1=0,\,a_{n+1}=a_n+\frac{1}{n(n+1)}$\\
            \bookterm{partial}{部分分数分解}を用いると,
            \[
              \frac{1}{k(k+1)}
              =\frac{1}{k}-\frac{1}{k+1}
            \]
            であるから,
            \begin{align*}
              a_n
                &=\sum_{k=1}^{n-1}\frac{1}{k(k+1)}\\
                &=\sum_{k=1}^{n-1}
                  \left(\frac{1}{k}-\frac{1}{k+1}\right)\\
                &=\left(1-\frac12\right)
                  +\left(\frac12-\frac13\right)\\
                &\quad +\cdots+\left(\frac{1}{n-1}-\frac1n\right)\\
                &=1-\frac1n
            \end{align*}

      \item $\displaystyle a_1=1,\,a_{n+1}=a_n+(n+1)^2$\\
            \bookterm{difference}{階差}型の\bookterm{general}{一般項}を用いると,
            \begin{align*}
              a_n
                &=1+\sum_{k=1}^{n-1}(k+1)^2\\
                &=1+\sum_{j=2}^{n}j^2\\
                &=\sum_{j=1}^{n}j^2\\
                &=\frac{1}{6}n(n+1)(2n+1)
            \end{align*}
    \end{enumerate}
